\documentclass{report}
\usepackage{amsmath,amssymb}
\setlength{\parindent}{0mm}
\setlength{\parskip}{1em}
\begin{document}
\begin{center}
\rule{6in}{1pt} \
{\large James Adler \\
{\bf Numerical Approximation of Asymptotically Disappearing Solutions of Maxwell�s Equations}}

503 Boston Ave \\ Tufts University \\ Mathematics Department \\ Medford \\ MA 02155
\\
{\tt james.adler@tufts.edu}\\
Vesselin  Petkov\\
Ludmil Zikatanov\end{center}

This work is on the numerical approximation of incoming solutions to
Maxwell's equations, whose energy decays exponentially with time
(asymptotically disappearing), meaning that the leading term of the
back-scattering matrix becomes negligible. For the exterior of a sphere,
such solutions are obtained by Colombini, Petkov and Rauch by specifying
a maximal dissipative boundary condition on the sphere and setting
appropriate initial conditions.

We consider a mixed finite element approximation of Maxwell's equations
in the exterior of a polyhedron whose boundary approximates the sphere.
We use the standard Nedelec-Raviart-Thomas elements and a Crank-Nicholson
scheme to approximate the electric and magnetic fields. We set discrete
initial conditions with standard interpolation, modified so that these
initial conditions are divergence-free. We prove that with such initial
conditions, the fully discrete approximation to the electric field is
weakly divergence-free for all time. We show numerically that the
finite-element solution approximates well the asymptotically disappearing
solutions constructed analytically when the mesh size becomes small.


\end{document}
